% Generado con herramientas/estadisticas_etl.py; tiempos de Hop fijados a la ejecución de la Figura 4.2 (18-11-2026 21:26).
\begin{table}[!htbp]
  \centering
  \caption{Tiempo de ejecución por etapa}
  \label{tab:est_tiempos}
  \small
  \begin{tabularx}{\linewidth}{@{}>{\raggedright\arraybackslash}Xr@{}}
    \toprule
    \textbf{Etapa} & \textbf{Segundos} \\
    \midrule
    Preparación en Python: Buses InterUrbanos & 4,4 \\
    Preparación en Python: Trenes & 0,1 \\
    Preparación en Python: Radio Taxis & 0,4 \\
    Preparación en Python: Logística y Distribución & 1,4 \\
    Preparación en Python: Aerolíneas & 90,6 \\
    Preparación en Python: URapid Delivery & 0,9 \\
    \midrule
    \textit{Subtotal preparación en Python} & \textit{97,8} \\
    \midrule
    Hop: pipeline 01\_staging & 3,1 \\
    Hop: pipeline 02\_dimensiones & 0,3 \\
    Hop: pipeline 03\_hechos & 10,6 \\
    \midrule
    \textit{Subtotal workflow principal de Hop} & \textit{15,0} \\
    \midrule
    \textbf{Total del ETL} & \textbf{112,8} \\
    \bottomrule
  \end{tabularx}
  \par\smallskip\footnotesize El subtotal de Hop es la duración del workflow principal: además de los tres pipelines, ejecuta los scripts SQL de staging, de preparación de la carga y de validación. El total suma la preparación de las seis líneas y el workflow. Los tiempos de Python provienen del registro de la preparación; los de Hop, de la ejecución de la Figura~\ref{fig:etl_terminal}.
\end{table}
